% !TeX root = ../sections/02_exercises.tex

\begin{exercise}
	次の行列 \(A\) の単因子 \(a_1,\ldots,a_r\) を求めよ。
	また，
	\[
	PAQ=\operatorname{diag}(a_1,\ldots,a_r)
	\]
	が成り立つようなユニモジュラー行列 \(P,Q\) を求めよ。
	
	\begin{enumerate}
		\item
		\[
		A=
		\begin{pmatrix}
			2 & 0 & 0\\
			0 & 3 & 0\\
			0 & 0 & 4
		\end{pmatrix}
		\]
		
		\item
		\[
		A=
		\begin{pmatrix}
			2 & 1\\
			-4 & 2
		\end{pmatrix}
		\]
		
		\item
		\[
		A=
		\begin{pmatrix}
			2 & 1 & -1\\
			-2 & -3 & 1\\
			-2 & -1 & 3
		\end{pmatrix}
		\]
	\end{enumerate}
\end{exercise}

\begin{background}
	整数行列 \(A\) に対して，整数係数の基本行変形と基本列変形を行うことで，
	対角行列
	\[
	\operatorname{diag}(a_1,\ldots,a_r)
	\]
	に変形できる。
	ただし，各 \(a_i\) は正の整数として選び，
	\[
	a_1\mid a_2\mid \cdots \mid a_r
	\]
	を満たすようにする。
	この \(a_1,\ldots,a_r\) を \(A\) の単因子という。
	
	整数係数の基本行変形は左からユニモジュラー行列を掛けることに対応し，
	整数係数の基本列変形は右からユニモジュラー行列を掛けることに対応する。
	したがって，行列 \(A\) を単因子標準形に変形する過程を記録すれば，
	\[
	PAQ=\operatorname{diag}(a_1,\ldots,a_r)
	\]
	となるユニモジュラー行列 \(P,Q\) も得られる。
	
	また，単因子を確認するためには，小行列式の最大公約数を用いる方法がある。
	\[
	\Delta_k
	=
	\gcd\{\text{\(A\) の \(k\) 次小行列式}\}
	\]
	とおくと，
	\[
	a_1=\Delta_1,
	\qquad
	a_1a_2=\Delta_2,
	\qquad
	a_1a_2a_3=\Delta_3
	\]
	が成り立つ。
\end{background}

\begin{strategy}
	この問題では，まず単因子を求め，その後で実際の基本変形を追って
	ユニモジュラー行列 \(P,Q\) を求める。
	
	\begin{enumerate}
		\item[(1)]
		行列はすでに対角行列であるが，
		対角成分 \(2,3,4\) は
		\[
		2\mid 3\mid 4
		\]
		を満たしていない。
		そのため，整数係数の行基本変形・列基本変形を用いて，
		互いに素な成分を分解しながら
		\[
		a_1\mid a_2\mid a_3
		\]
		となる形に直す。
		確認には
		\[
		\Delta_1,\quad \Delta_2,\quad \Delta_3
		\]
		を計算するとよい。
		
		\item[(2)]
		まず全成分の最大公約数を計算して \(\Delta_1\) を求める。
		次に，行列式
		\[
		\det A
		\]
		を計算して \(\Delta_2\) を求める。
		すると
		\[
		a_1=\Delta_1,
		\qquad
		a_2=\frac{\Delta_2}{\Delta_1}
		\]
		で単因子が決まる。
		その後，基本行変形・基本列変形によって実際に
		\(\operatorname{diag}(a_1,a_2)\) に変形し，
		その操作を \(P,Q\) に反映させる。
		
		\item[(3)]
		まず
		\[
		\Delta_1
		\]
		を全成分の最大公約数として求める。
		次に，すべての \(2\) 次小行列式の最大公約数から
		\[
		\Delta_2
		\]
		を求める。
		最後に，
		\[
		\Delta_3=|\det A|
		\]
		を計算する。
		これにより
		\[
		a_1=\Delta_1,
		\qquad
		a_2=\frac{\Delta_2}{\Delta_1},
		\qquad
		a_3=\frac{\Delta_3}{\Delta_2}
		\]
		として単因子が求まる。
		
		実際に \(P,Q\) を求めるには，
		まず絶対値が小さい非零成分を左上に移動させ，
		ユークリッドの互除法に対応する行変形・列変形を繰り返して，
		左上成分が第1行・第1列のすべての成分を割るようにする。
		その後，第1行・第1列を消去し，残った部分行列に同じ操作を繰り返す。
	\end{enumerate}
\end{strategy}